\subsection{Basic Strategy for Larger Systems} 
Let's consider our basic strategy for solving a larger system.
\begin{itemize}
\item When we had a system of 2 equations with 2 variables we \makeblank[1in] them to get \makeblanksmall equation with \makeblanksmall variable.
\item We could then \makeblank[3in].
\item In general, when we have 3 equations with 3 variables, we will combine them to give \makeblanksmall equations with \makeblanksmall variables.
\item We then solve this by \makeblank.
\item Note that if we have 3 equations, there are \makeblanksmall possible combinations. We will choose \makeblanksmall of them.
\item If we have \makeblanksmall equations, there are \makeblanksmall possible combinations! We will only choose \makeblanksmall of them.
\end{itemize}
The hardest part of solving a larger system is \makeblank of what we're doing. In order to solve such a system, we will \makeblank[1in] the equations and \makeblank[1in] what we're doing at each step.